\section{Quando o modelo de aprendizagem profunda falha?}
\label{sec:falhas}
\subsection{A CNN reage à troca e inverte o sentido do contraste}
A termina em \textit{wonderful} (positiva) e B em \textit{terrible} (negativa).
Com três sementes por tokenizador, os seis treinos previram negativa em A e
positiva em B. O escore positivo foi de \PalavrasA{} a \PalavrasB{} com palavras
e de \BPEA{} a \BPEB{} com BPE, na média.
\href{https://costaleirbag.github.io/projeto-mae5920/\#slide-9}{Troca interativa A/B}.

\subsection{O máximo global apaga a posição entre janelas}
Uma frase repetida separa \textit{I loved the acting} de \textit{I hated the ending}.
Trocamos os dois blocos e verificamos que as janelas de 3, 4 e 5 IDs formam
os mesmos multiconjuntos nos dois tokenizadores. Cada filtro recebe as mesmas
janelas e conserva o mesmo máximo. Os logits coincidiram nos seis treinos.

Com os filtros congelados, guardamos dois picos por filtro. Uma cabeça linear
recebe menor e maior pico. A outra recebe os picos na ordem do texto. Guardar
a posição fez os escores de A e B diferirem. Em \NParesHumanos{} pares de
resenhas editadas, o acerto nas duas versões foi de \SemPosHumanos{} sem
posição e \ComPosHumanos{} com posição. \fonteatual{seminario/diagnosticos.md}{Método e três sementes}.

\subsection{O sentimento fora do IMDb}
Avaliamos os mesmos pesos em 1.000 frases de Amazon e 1.000 de Yelp, com
500 exemplos por classe em cada fonte \cite{kotzias}. Gênero e comprimento
mudam juntos em relação ao IMDb.

\begin{tabularx}{\linewidth}{@{}Xrrrrrr@{}}
\toprule
 & \multicolumn{3}{c}{\tabelahead{Palavras}} & \multicolumn{3}{c}{\tabelahead{BPE}}\\
\cmidrule(lr){2-4}\cmidrule(lr){5-7}
\tabelahead{Fonte} & \tabelahead{Acurácia} & \tabelahead{F1 macro} & \tabelahead{CE} &
\tabelahead{Acurácia} & \tabelahead{F1 macro} & \tabelahead{CE}\\
\midrule
\LinhasForaDominio
\bottomrule
\end{tabularx}
\legendinha{Tabela 11. Médias de três sementes em frases selecionadas por
polaridade clara. \fonteatual{seminario/dados/fora_dominio_resultados.json}{Resultados por semente}.}

Em \textit{The Room}, todos os treinos leram o elogio seguido de “No one”
como positivo. Os comentários de \textit{Baby} deram positivo. O de
\textit{Friday} deu negativo. \fonteatual{seminario/dados/fora_dominio_resultados.json}{Cinco casos sem rótulo}.

\subsection{Tom e fatos na geração do Qwen}
No bloco 28, somamos \(\alpha s v\) após \texttt{REVIEW}, com \(\alpha=-2,0,+2\).
A direção \(v\) compara as ativações médias para dois tons em 128 enredos.
O Qwen, como juiz automático, reconheceu os dois tons-alvo em 55 de 64 fontes.
\texttt{FACT} ficou igual. Em um caso, \texttt{REVIEW} inventou um cenário de
cidade pequena. \fonte{results/qwen_improved/fact_flag_reading.json}{Alertas factuais}.

\begin{thebibliography}{9}
\small
\setlength{\itemsep}{0pt}
\setlength{\parskip}{0pt}
\bibitem{maas} Maas, A. L. et al. \textit{Learning Word Vectors for Sentiment Analysis}.
ACL, 2011. \href{https://aclanthology.org/P11-1015/}{ACL Anthology}.
\bibitem{kim} Kim, Y. \textit{Convolutional Neural Networks for Sentence Classification}.
EMNLP, 2014. \href{https://aclanthology.org/D14-1181/}{ACL Anthology}.
\bibitem{sennrich} Sennrich, R., Haddow, B. e Birch, A.
\textit{Neural Machine Translation of Rare Words with Subword Units}. ACL, 2016.
\href{https://aclanthology.org/P16-1162/}{ACL Anthology}.
\bibitem{prince} Prince, S. J. D. \textit{Understanding Deep Learning}.
MIT Press, 2023. Caps. 6, 8, 9 e 10. \href{https://udlbook.github.io/udlbook/}{Livro e materiais}.
\bibitem{devlin} Devlin, J. et al. \textit{BERT: Pre-training of Deep Bidirectional
Transformers for Language Understanding}. NAACL, 2019.
\href{https://aclanthology.org/N19-1423/}{ACL Anthology}.
\bibitem{qwen} Yang, A. et al. \textit{Qwen3 Technical Report}, 2025.
\href{https://arxiv.org/abs/2505.09388}{arXiv:2505.09388}.
\bibitem{kotzias} Kotzias, D. et al. \textit{From Group to Individual Labels Using Deep Features}.
KDD, 2015. \href{https://doi.org/10.1145/2783258.2783380}{Artigo}.
\href{https://doi.org/10.24432/C57604}{Coleção UCI}.
\end{thebibliography}
